\section{The integral two-term block in the high-support class}\label{sec:high}

Assume throughout this section that $g=by^{2p}+cy^\ell+R_1(y)$,
where $\deg R_1<\ell$, $p+2\le\ell<2p$, and $b,c\ne0$.
No multiple of $p$ lies strictly between $\ell$ and $2p$.
With $h=2p-\ell$, the hypotheses $q\ge p^3$ and $1\le h\le p-2$
imply the two strict inequalities
\begin{equation}\label{eq:strict-gaps}
 q-2p>h,\qquad
 q(p-h-1)-2p(p-1)>h.
\end{equation}
For the second, $p-h-1\ge1$ and
\[
 q(p-h-1)-2p(p-1)
 \ge p^3-2p(p-1)>p-2\ge h.
\]
These margins are uniform over the entire high-support coefficient
class, including $q=p^3$.

\begin{lemma}[First nonzero binomial index]\label{lem:binomial}
If $D\ge1$ and $r_D=p^{v_p(D)}$, then
\begin{equation}\label{eq:binomial-bound}
 \deg\bigl((y^q+u_0)^D-y^{qD}\bigr)
 \le q(D-r_D)+2p r_D.
\end{equation}
The first potentially nonzero binomial index is $r_D$.
\end{lemma}
\begin{proof}
Write $D=r_D d$ with $p\nmid d$. The identity
$(Y+Z)^D=(Y^{r_D}+Z^{r_D})^d$ shows that the first nonconstant
term is $dY^{r_D(d-1)}Z^{r_D}$. Since $\deg u_0=2p<q$, a term
with a larger index has strictly smaller total degree than the
one with index $r_D$. This proves the bound.
\end{proof}

\begin{proposition}[High-support invariant]\label{prop:high}
Define integers and nonzero coefficients by
\[
 A_1=0,\qquad A_{j+1}=q(A_j+\ell-1),\qquad
 C_j=(c\bar\ell)^{j-1}.
\]
Then for every $j\ge1$,
\begin{equation}\label{eq:high-invariant}
 \delta^j y=C_jy^{A_j}(by^{2p}+cy^\ell)+R_j,
 \qquad \deg R_j<A_j+\ell.
\end{equation}
In particular \eqref{eq:leading-claim} holds with the high-support
degree in \eqref{eq:degree-formulas}.
\end{proposition}

\begin{proof}
At $j=1$, the assertion is $\delta y=g(y)-ax$; the omitted terms
have degree below $\ell$. Suppose \eqref{eq:high-invariant} holds
at $j$, and set
\[
 D=A_j+2p,\qquad E=A_j+\ell,\qquad
 d=q(E-1)+2p.
\]
The integer $A_j$ is zero or divisible by $q$. Because $2p<q$,
$v_p(D)=1$, whereas $E\equiv\ell\pmod p$ is nonzero.

Consider first $C_jc\delta(y^E)$. Its binomial-index-one term is
\[
 C_jc\bar\ell\,y^{q(E-1)}(g(y)-ax).
\]
This supplies precisely the two displayed terms at degrees $d$
and $d-h$, with common multiplier $C_{j+1}$. Every other term
in this index-one contribution has smaller degree than $d-h$.
Each term with index $i\ge2$ has degree at most
\[
 q(E-i)+2pi=d-(i-1)(q-2p)<d-h
\]
by the first inequality in \eqref{eq:strict-gaps}.

For the leading monomial $C_jb\delta(y^D)$, Lemma~\ref{lem:binomial}
gives degree at most $q(D-p)+2p^2$. Its deficit from $d$ is
\[
 d-\bigl(q(D-p)+2p^2\bigr)
 =q(p-h-1)-2p(p-1)>h.
\]
Thus this term changes neither retained coefficient. Finally,
\eqref{eq:ordinary-degree-bound} and the integral degree of $R_j$
give
\[
 \deg\delta R_j\le q(E-1)=d-2p<d-h.
\]
All omitted terms lie below the second retained term. This proves
\eqref{eq:high-invariant} at $j+1$.

The top degree is $D_j=A_j+2p$, and hence
\[
 D_1=2p,\qquad D_{j+1}=q(D_j-h-1)+2p.
\]
Solving this affine recurrence yields the first formula in
\eqref{eq:degree-formulas}. Its coefficient is $bC_j$, as in
\eqref{eq:leading-claim}. Positivity and integrality follow from
$A_j\ge0$. Also $D_1<q$ and $D_{j+1}<qD_j$, since
$q(h+1)>2p$, so $D_j<q^j$ for every $j$.
\end{proof}

The proof uses strict margins to prevent a lower term from entering
the retained block. Those margins no longer apply at the low-support
boundary $\ell=p+1$. A different state is needed there, including
when the coefficient of $y^{p+1}$ vanishes altogether.
